\scanpage{35}
\begin{center}- 187 -\end{center}
\hfill n° 391

les systèmes de racines correspondant à ces schémas, ce qui montrera \uline{a fortiori} l'existence des groupes de Coxeter correspondants.

Soient $R$ un système de racines réduit de rang $\ell$, $S(C)$ un système de racines simples relatif à une chambre $C$ de $R$. Si $u,v$ sont deux éléments distincts de $S(C)$, on sait, (§ 6, n° 5), que $(u|v)\leq0$, et par suite l'entier $n(u,v)=2(u|v)/(u|u)$ ne peut prendre que les valeurs $0,-1,-2$ ou $-3$, (§ 6, n° 3) ; en outre, dans les trois derniers cas les valeurs de $\|v\|/\|u\|$ sont respectivement $1$, $\sqrt2$ et $\sqrt3$, et l'ordre de $s_us_v$ est 2 dans le premier cas, 3, 4 et 6 respectivement dans les trois derniers.

Nous appellerons \uline{schéma de Dynkin} un couple formé d'un graphe fini $X$ et d'une application $f$ de $X$ dans l'ensemble $\{1,2,3\}$. Etant donnés deux graphes de Dynkin $(X,f)$, $(X',f')$, un \uline{isomorphisme} de $(X,f)$ sur $(X',f')$ est une bijection $\varphi:X\longrightarrow X'$ qui est un isomorphisme des graphes sous-jacents et telle que $f'(\varphi(i))=f(i)$ pour tout $i\in X$. Etant donné un système de racines réduit $R$ et une chambre $C$ de $R$, on appelle \uline{schéma de Dynkin de} $R$ relatif à $C$ le schéma de Dynkin dont le graphe sous-jacent a pour sommets les racines simples de $S(C)$, les arêtes étant les parties $\{u,v\}$ à deux éléments de $S(C)$ telles que $(u|v)\ne0$, et l'application $f:X\longrightarrow\{1,2,3\}$ étant telle que $f$ soit égale à 1 dans l'ensemble des racines de $S(C)$ de plus petite longueur $\lambda$, à 2 dans l'ensemble des racines de $S(C)$ de longueur $\lambda\sqrt2$, à 3 dans l'ensemble des racines de $S(C)$ de longueur $\lambda\sqrt3$ ; on sait d'ailleurs qu'en fait, ou bien $f$ est constante et égale à 1 dans $X$, ou bien ne prend que les valeurs 1 et 2, ou bien ne prend que les valeurs 1 et 3, (§ 6, n° 4, prop. 6).
\mspage{34}
